% AI tự sinh cho Ôn thi HK Lớp 10

\begin{bt}%[0H5C2-6]%[ma:a1b2c3d4]
% nguon: AI tự sinh (Ôn thi HK)
Một chiếc thuyền xuất phát từ bờ $A$, cần sang bờ đối diện tại điểm $B$ nằm chính đối diện (theo phương vuông góc với bờ). Dòng nước chảy song song với bờ với vectơ vận tốc $\overrightarrow{v_n}$ có độ lớn $3$ km/h. Thuyền có vận tốc riêng (so với nước) là $\overrightarrow{v_t}$ có độ lớn $5$ km/h. Hỏi người lái phải hướng mũi thuyền lệch một góc $\alpha$ (so với phương vuông góc với bờ) bao nhiêu độ để thuyền đi đúng hướng $AB$? Biết vận tốc thực của thuyền là tổng $\overrightarrow{v} = \overrightarrow{v_t} + \overrightarrow{v_n}$. Kết quả làm tròn đến hàng đơn vị (độ).
\shortans{37}
	\loigiai{
Gọi $Oxy$ với trục $Ox$ song song dòng nước, $Oy$ vuông góc bờ hướng từ $A$ sang $B$.
Vectơ vận tốc nước: $\overrightarrow{v_n} = (3; 0)$.
Mũi thuyền lệch góc $\alpha$ so với $Oy$ về phía ngược dòng (để bù dòng chảy): $\overrightarrow{v_t} = (-5\sin\alpha;\, 5\cos\alpha)$.
Vận tốc thực: $\overrightarrow{v} = (3 - 5\sin\alpha;\, 5\cos\alpha)$.
Để thuyền đi đúng hướng $AB$ (theo $Oy$), thành phần $x$ phải bằng $0$:
$$3 - 5\sin\alpha = 0 \Rightarrow \sin\alpha = \frac{3}{5} \Rightarrow \alpha \approx 36{,}87^\circ.$$
Làm tròn đến hàng đơn vị: $\alpha \approx 37^\circ$.
}

\end{bt}

\begin{bt}%[0H5C2-2]%[ma:b2c3d4e5]
% nguon: AI tự sinh (Ôn thi HK)
Trong mặt phẳng tọa độ $Oxy$, cho tam giác $ABC$ với $A(1;3)$, $B(-2;1)$, $C(4;-1)$. Gọi $M$ là điểm thỏa mãn $\overrightarrow{MA} + 2\overrightarrow{MB} - \overrightarrow{MC} = \overrightarrow{0}$. Tìm tọa độ điểm $N$ trên đường thẳng $d: x + y - 2 = 0$ sao cho $MN$ vuông góc với $AB$. Tính $x_N + 2y_N$.
\shortans{12,5}
	\loigiai{
Từ $\overrightarrow{MA} + 2\overrightarrow{MB} - \overrightarrow{MC} = \overrightarrow{0}$, gọi $M(x;y)$:
$$\overrightarrow{MA}=(1-x;3-y),\ \overrightarrow{MB}=(-2-x;1-y),\ \overrightarrow{MC}=(4-x;-1-y).$$
Phương trình hoành độ: $(1-x)+2(-2-x)-(4-x)=0 \Rightarrow 1-x-4-2x-4+x=0 \Rightarrow -2x-7=0 \Rightarrow x=-\dfrac{7}{2}$.
Phương trình tung độ: $(3-y)+2(1-y)-(-1-y)=0 \Rightarrow 3-y+2-2y+1+y=0 \Rightarrow -2y+6=0 \Rightarrow y=3$.
Vậy $M\left(-\dfrac{7}{2};3\right)$.
Vectơ $\overrightarrow{AB}=(-3;-2)$. Điều kiện $MN \perp AB$: $\overrightarrow{MN}\cdot\overrightarrow{AB}=0$.
$N \in d \Rightarrow N(t;\,2-t)$. $\overrightarrow{MN}=\left(t+\dfrac{7}{2};\, -1-t\right)$.
$$\left(t+\dfrac{7}{2}\right)(-3) + (-1-t)(-2)=0 \Rightarrow -3t-\dfrac{21}{2}+2+2t=0 \Rightarrow -t-\dfrac{17}{2}=0 \Rightarrow t=-\dfrac{17}{2}.$$
$N\left(-\dfrac{17}{2};\, \dfrac{21}{2}\right)$. Vậy $x_N + 2y_N = -\dfrac{17}{2} + 21 = \dfrac{25}{2} = 12{,}5$.
}

\end{bt}

\begin{bt}%[0H5C2-5]%[ma:c3d4e5f6]
% nguon: AI tự sinh (Ôn thi HK)
Cho tam giác $ABC$ có $AB=4$, $AC=6$, góc $\widehat{BAC}=60^\circ$. Gọi $M$ là điểm di động trên đường thẳng $BC$. Tìm giá trị nhỏ nhất của biểu thức
$$P = \left|\overrightarrow{MA}+\overrightarrow{MB}\right| + \left|\overrightarrow{MA}+\overrightarrow{MC}\right|.$$
Kết quả làm tròn đến hàng phần mười.
\shortans{5,3}
	\loigiai{
Gọi $I$, $J$ lần lượt là trung điểm $AB$, $AC$. Khi đó $\overrightarrow{MA}+\overrightarrow{MB}=2\overrightarrow{MI}$, $\overrightarrow{MA}+\overrightarrow{MC}=2\overrightarrow{MJ}$.
Suy ra $P = 2(MI + MJ)$.
Ta có $MI + MJ \ge IJ$ với mọi $M$, dấu bằng khi $M$ nằm trên đoạn $IJ$ và $M \in BC$, tức $M = IJ \cap BC$.
Tính $IJ$: $IJ$ là đường trung bình của tam giác $ABC$ ứng với cạnh $BC$, nên $IJ = \dfrac{BC}{2}$.
Tính $BC$: $BC^2 = AB^2+AC^2-2\cdot AB\cdot AC\cos60^\circ = 16+36-24=28 \Rightarrow BC=2\sqrt{7}$.
Vậy $IJ = \sqrt{7}$, và $P_{\min} = 2\sqrt{7} \approx 5{,}3$.
}

\end{bt}

\begin{bt}%[0H5C2-3]%[ma:d4e5f6a7]
% nguon: AI tự sinh (Ôn thi HK)
Trong mặt phẳng tọa độ $Oxy$, cho hai điểm $A(2;1)$, $B(-1;3)$ và đường thẳng $\Delta: x - 2y + 4 = 0$. Biết điểm $M \in \Delta$ thỏa mãn $\left|\overrightarrow{MA}+\overrightarrow{MB}\right|$ đạt giá trị nhỏ nhất. Gọi $M_0$ là điểm đó. Tính $5x_{M_0} - y_{M_0}$.
\shortans{-0,2}
	\loigiai{
Gọi $I$ là trung điểm $AB$: $I\left(\dfrac{1}{2};2\right)$. Ta có $\overrightarrow{MA}+\overrightarrow{MB}=2\overrightarrow{MI}$, nên $\left|\overrightarrow{MA}+\overrightarrow{MB}\right|=2MI$ nhỏ nhất khi $MI$ nhỏ nhất, tức $M$ là hình chiếu vuông góc của $I$ trên $\Delta$.
Đường thẳng qua $I$ vuông góc $\Delta$ có vectơ pháp tuyến $\overrightarrow{n}=(1;-2)$ là vectơ chỉ phương, nên nhận $\overrightarrow{u}=(2;1)$ làm pháp tuyến:
$$2\left(x-\dfrac{1}{2}\right)+1(y-2)=0 \Rightarrow 2x+y-3=0.$$
Tọa độ $M_0$ là nghiệm hệ:
$$\begin{cases} x-2y+4=0 \\ 2x+y-3=0 \end{cases} \Rightarrow \begin{cases} x=2y-4 \\ 2(2y-4)+y-3=0 \end{cases} \Rightarrow 5y=11 \Rightarrow y=\dfrac{11}{5},\ x=\dfrac{2}{5}.$$
Vậy $5x_{M_0}-y_{M_0} = 5\cdot\dfrac{2}{5}-\dfrac{11}{5} = 2-\dfrac{11}{5} = -\dfrac{1}{5} = -0{,}2$.
}

\end{bt}

\begin{bt}%[0H5C2-4]%[ma:e5f6a7b8]
% nguon: AI tự sinh (Ôn thi HK)
Trong mặt phẳng tọa độ $Oxy$, cho tam giác $ABC$ có $A(1;2)$, $B(3;-1)$, $C(-1;4)$. Biết điểm $M$ thỏa mãn $\overrightarrow{MA}+\overrightarrow{MB}+\overrightarrow{MC}=\overrightarrow{0}$ và điểm $N$ thỏa mãn $\overrightarrow{NA}-2\overrightarrow{NB}+\overrightarrow{NC}=\overrightarrow{0}$. Tính độ dài đoạn thẳng $MN$ (làm tròn đến hàng phần trăm).
\shortans{0,83}
	\loigiai{
\textbf{Bước 1: Tìm $M$ (trọng tâm).}
$M$ là trọng tâm tam giác: $M\left(\dfrac{1+3-1}{3};\dfrac{2-1+4}{3}\right) = (1;\,5/3)$.

\textbf{Bước 2: Tìm $N$.}
$\overrightarrow{NA}-2\overrightarrow{NB}+\overrightarrow{NC}=\overrightarrow{0}$. Gọi $N(x;y)$:
$$\overrightarrow{NA}=(1-x;2-y),\ \overrightarrow{NB}=(3-x;-1-y),\ \overrightarrow{NC}=(-1-x;4-y).$$
Hoành độ: $(1-x)-2(3-x)+(-1-x)=0 \Rightarrow 1-x-6+2x-1-x=0 \Rightarrow -6=0$ (vô lý).

Nhận xét: hệ số $1-2+1=0$ nên đẳng thức không xác định điểm $N$ duy nhất mà tương đương điều kiện vectơ. Thật vậy:
$$\overrightarrow{NA}-2\overrightarrow{NB}+\overrightarrow{NC} = (\overrightarrow{NA}+\overrightarrow{NC})-2\overrightarrow{NB} = 2\overrightarrow{NM'}-2\overrightarrow{NB}$$
với $M'$ trung điểm $AC$: $M'\left(0;3\right)$. Điều kiện trở thành $\overrightarrow{NM'}=\overrightarrow{NB}$, tức $N$ là trung điểm $BM'$:
$$N\left(\dfrac{3+0}{2};\dfrac{-1+3}{2}\right) = \left(\dfrac{3}{2};1\right).$$

\textbf{Bước 3: Tính $MN$.}
$$MN = \sqrt{\left(\dfrac{3}{2}-1\right)^2 + \left(1-\dfrac{5}{3}\right)^2} = \sqrt{\dfrac{1}{4}+\dfrac{4}{9}} = \sqrt{\dfrac{25}{36}} = \dfrac{5}{6} \approx 0{,}83.$$
}

\end{bt}

\begin{bt}%[0H5C2-6]%[ma:f6a7b8c9]
% nguon: AI tự sinh (Ôn thi HK)
Một người kéo một vật nặng trên mặt phẳng nằm ngang bằng hai sợi dây buộc vào cùng một điểm trên vật. Lực kéo trên dây thứ nhất có độ lớn $F_1 = 80$ N, dây thứ hai có độ lớn $F_2 = 60$ N. Hai dây hợp với nhau một góc $60^\circ$. Hợp lực $\overrightarrow{F} = \overrightarrow{F_1}+\overrightarrow{F_2}$ có độ lớn bằng bao nhiêu N? (làm tròn đến hàng đơn vị).
\shortans{122}
	\loigiai{
Bình phương hai vế $\overrightarrow{F}=\overrightarrow{F_1}+\overrightarrow{F_2}$:
$$F^2 = F_1^2 + F_2^2 + 2F_1F_2\cos60^\circ = 80^2 + 60^2 + 2\cdot 80\cdot 60\cdot\dfrac{1}{2}.$$
$$F^2 = 6400 + 3600 + 4800 = 14800 \Rightarrow F = \sqrt{14800} = 20\sqrt{37} \approx 121{,}66\ \text{N}.$$
Làm tròn đến hàng đơn vị: $F \approx 122$ N.
}

\end{bt}

\begin{bt}%[0H5C2-1]%[ma:a7b8c9d0]
% nguon: AI tự sinh (Ôn thi HK)
Trong mặt phẳng tọa độ $Oxy$, cho tam giác $ABC$ có $A(-1;2)$, $B(3;4)$, $C(0;-3)$. Điểm $M$ thỏa mãn $\overrightarrow{AM} = 2\overrightarrow{AB} - 3\overrightarrow{AC}$. Điểm $N$ đối xứng với $M$ qua trung điểm $I$ của $BC$. Tìm tung độ của $N$.
\shortans{-20}
	\loigiai{
$\overrightarrow{AB}=(4;2)$, $\overrightarrow{AC}=(1;-5)$.
$$\overrightarrow{AM} = 2(4;2)-3(1;-5) = (8;4)-(3;-15) = (5;19).$$
Suy ra $M = A + \overrightarrow{AM} = (-1+5;\,2+19) = (4;21)$.
Trung điểm $I$ của $BC$: $I\left(\dfrac{3+0}{2};\dfrac{4-3}{2}\right) = \left(\dfrac{3}{2};\dfrac{1}{2}\right)$.
$N$ đối xứng với $M$ qua $I$ nên $I$ là trung điểm $MN$:
$$x_N = 2\cdot\dfrac{3}{2}-4 = -1,\quad y_N = 2\cdot\dfrac{1}{2}-21 = 1-21 = -20.$$
Vậy tung độ của $N$ là $-20$.
}

\end{bt}

\begin{bt}%[0H5C2-2]%[ma:b8c9d0e1]
% nguon: AI tự sinh (Ôn thi HK)
Cho tam giác $ABC$ đều cạnh $a$. Điểm $M$ thay đổi trên đường tròn ngoại tiếp tam giác $ABC$. Tìm giá trị lớn nhất của biểu thức
$$T = \left|\overrightarrow{MA}+\overrightarrow{MB}\right| + \left|\overrightarrow{MB}+\overrightarrow{MC}\right| + \left|\overrightarrow{MC}+\overrightarrow{MA}\right|.$$
Kết quả làm tròn đến hàng phần trăm theo $a$ (tức tính hệ số $k$ với $T_{\max}=k\cdot a$).
\shortans{4,56}
	\loigiai{
Gọi $I$, $J$, $K$ lần lượt là trung điểm $AB$, $BC$, $CA$. Khi đó:
$$T = 2(MI + MJ + MK).$$
Với $M$ trên đường tròn ngoại tiếp tâm $O$ bán kính $R = \dfrac{a}{\sqrt{3}}$. Các điểm $I,J,K$ là trung điểm các cạnh, cách $O$ một khoảng $r = \dfrac{R}{2} = \dfrac{a}{2\sqrt{3}}$ (đường tròn Euler).
Với mỗi điểm $P$ cố định, $MP \le MO + OP = R + OP$. Do đó:
$$MI \le R + OI,\quad MJ \le R + OJ,\quad MK \le R + OK.$$
Dấu bằng đồng thời xảy ra khi $M$ nằm trên tia đối của $OI$, $OJ$, $OK$ — không thể đồng thời. Do đó dùng đánh giá tổng quát chặt hơn: $MI+MJ+MK \le 3R + (OI+OJ+OK)$ chỉ là cận trên lỏng.

Xét chính xác: $T = 2(MI+MJ+MK)$. Vì $I,J,K$ tạo thành tam giác đều tâm $O$ bán kính $r=\dfrac{a}{2\sqrt{3}}$, và $M$ trên đường tròn tâm $O$ bán kính $R=2r$. Khi $M$ trùng với một đỉnh "xa" của đường tròn lớn (ví dụ $M$ đối xứng với $I$ qua $O$), thì $MI = R+r = 3r$, còn $MJ=MK$ tính được:
$$MJ^2 = R^2+r^2-2Rr\cos120^\circ = 4r^2+r^2+r^2 = 6r^2 \Rightarrow MJ = r\sqrt{6}.$$
Kiểm tra tổng: $MI+MJ+MK = 3r + 2r\sqrt{6} = r(3+2\sqrt{6})$.
Thử với $M$ là một trong ba điểm "gần" (trùng hướng $I$): $MI = R - r = r$, $MJ=MK=r\sqrt{6}$, tổng $= r(1+2\sqrt{6})$, nhỏ hơn. Vậy max đạt tại $M$ đối diện $I$: $MI+MJ+MK = r(3+2\sqrt{6})$.
$$T_{\max} = 2r(3+2\sqrt{6}) = 2\cdot\dfrac{a}{2\sqrt{3}}(3+2\sqrt{6}) = \dfrac{a}{\sqrt{3}}(3+2\sqrt{6}).$$
Hệ số $k = \dfrac{3+2\sqrt{6}}{\sqrt{3}} = \sqrt{3} + 2\sqrt{2} \approx 1{,}732+2{,}828 = 4{,}560 \approx 4{,}56$.
}

\end{bt}

\begin{bt}%[0H5C2-3]%[ma:c9d0e1f2]
% nguon: AI tự sinh (Ôn thi HK)
Cho tam giác $ABC$ có $AB=3$, $AC=4$, $BC=5$. Điểm $M$ di động sao cho $\left|\overrightarrow{MA}+\overrightarrow{MB}+\overrightarrow{MC}\right| = 2$. Gọi $G$ là trọng tâm tam giác. Tìm giá trị lớn nhất của biểu thức $Q = MA^2 + MB^2 + MC^2$. Kết quả làm tròn đến hàng phần mười.
\shortans{18}
	\loigiai{
Tam giác $ABC$ vuông tại $A$ (vì $3^2+4^2=5^2$). Trọng tâm $G$: $AG^2 = \dfrac{2AB^2+2AC^2-BC^2}{9} = \dfrac{18+32-25}{9} = \dfrac{25}{9}$, nên $AG=\dfrac{5}{3}$.

Điều kiện: $\left|\overrightarrow{MA}+\overrightarrow{MB}+\overrightarrow{MC}\right| = 3MG = 2 \Rightarrow MG = \dfrac{2}{3}$. Vậy $M$ chạy trên đường tròn tâm $G$ bán kính $\dfrac{2}{3}$.

Công thức: $MA^2+MB^2+MC^2 = 3MG^2 + GA^2+GB^2+GC^2$.
Tính $GA^2+GB^2+GC^2$: với $G$ trọng tâm, $GA^2+GB^2+GC^2 = \dfrac{AB^2+BC^2+CA^2}{3} = \dfrac{9+25+16}{3} = \dfrac{50}{3}$.
Và $3MG^2 = 3\cdot\dfrac{4}{9} = \dfrac{4}{3}$.
Vậy $Q = \dfrac{4}{3} + \dfrac{50}{3} = \dfrac{54}{3} = 18$ — không đổi với mọi $M$ thỏa điều kiện!

Do $Q$ không đổi, giá trị lớn nhất bằng $18$.
}

\end{bt}

\begin{bt}%[0H5C2-4]%[ma:d0e1f2a3]
% nguon: AI tự sinh (Ôn thi HK)
Trong mặt phẳng tọa độ $Oxy$, một vật chuyển động thẳng đều từ $A(1;2)$ đến $B(7;10)$ với vận tốc có độ lớn $5$ (đơn vị/giây). Cùng lúc đó một vật khác chuyển động từ $C(-2;6)$ theo hướng vectơ $\overrightarrow{u}=(1;-1)$ với vận tốc $3\sqrt{2}$ (đơn vị/giây). Biết rằng sau $t$ giây, vị trí hai vật lần lượt là $P$ và $Q$. Gọi $D$ là điểm sao cho $\overrightarrow{OD} = \overrightarrow{OP}+\overrightarrow{OQ}$ (tổng vectơ vị trí). Tìm quãng đường điểm $D$ đã di chuyển trong $2$ giây đầu tiên (làm tròn đến hàng phần mười).
\shortans{12,2}
	\loigiai{
\textbf{Bước 1: Vectơ vận tốc vật 1.}
$\overrightarrow{AB}=(6;8)$, $|\overrightarrow{AB}|=10$. Đơn vị hướng: $\left(\dfrac{3}{5};\dfrac{4}{5}\right)$. Vận tốc: $\overrightarrow{v_1}=5\cdot\left(\dfrac{3}{5};\dfrac{4}{5}\right)=(3;4)$.

\textbf{Bước 2: Vectơ vận tốc vật 2.}
Hướng $\overrightarrow{u}=(1;-1)$, $|\overrightarrow{u}|=\sqrt{2}$, đơn vị $\left(\dfrac{1}{\sqrt2};-\dfrac{1}{\sqrt2}\right)$. Vận tốc: $\overrightarrow{v_2}=3\sqrt2\cdot\left(\dfrac{1}{\sqrt2};-\dfrac{1}{\sqrt2}\right)=(3;-3)$.

\textbf{Bước 3: Vị trí $P$, $Q$ theo $t$.}
$P = A + t\overrightarrow{v_1} = (1+3t;\,2+4t)$.
$Q = C + t\overrightarrow{v_2} = (-2+3t;\,6-3t)$.

\textbf{Bước 4: Điểm $D$.}
$\overrightarrow{OD}=\overrightarrow{OP}+\overrightarrow{OQ} = (1+3t-2+3t;\, 2+4t+6-3t) = (-1+6t;\,8+t)$.
Vậy $D$ chuyển động thẳng đều với vận tốc $\overrightarrow{v_D}=(6;1)$, độ lớn $|\overrightarrow{v_D}|=\sqrt{37}$.

\textbf{Bước 5: Quãng đường trong $2$ giây.}
$s = |\overrightarrow{v_D}|\cdot 2 = 2\sqrt{37} \approx 2\cdot 6{,}083 = 12{,}166 \approx 12{,}2$.
}

\end{bt}
